\section{Discussion}

\label{sec:discussion}
% ══════════════════════════════════════════════════════════════════════════════

\subsection{Why Taxonomy Matters}

The success of taxonomic prototypes reflects a biological reality:
species within the same genus share morphological features due to common
ancestry. By explicitly modeling this hierarchical structure, TaxoNet
leverages phylogenetically informative features even when species-level
examples are scarce. The rank-attention mechanism acts as a soft decision
tree, routing queries to the appropriate level of taxonomic resolution
based on visual evidence quality.

\subsection{Limitations}

\paragraph{Cryptic species.} Species complexes that are
morphologically indistinguishable cannot be resolved by visual features
alone. Integration with molecular data (eDNA) or acoustic features is
needed.

\paragraph{Phenological variation.} Some species change appearance
dramatically across seasons. Explicit temporal prototype modulation
would be more principled than our current update mechanism.

\paragraph{Geographic bias.} Pre-training data is biased toward North
America and Europe. Performance on tropical and Southern Hemisphere
species may be lower due to representation gaps.

\subsection{Ethical Considerations}

\begin{itemize}
  \item \textbf{Poaching risk}: We implement geographic fuzzing for
    threatened species locations.
  \item \textbf{Over-reliance}: We display confidence calibration and
    encourage expert confirmation for low-confidence predictions.
  \item \textbf{Data sovereignty}: We follow the Nagoya Protocol and
    CARE principles for indigenous data governance.
\end{itemize}


% ══════════════════════════════════════════════════════════════════════════════
